\chapter*{Kata Pengantar}
Puji syukur penulis panjatkan kepada Tuhan Yang Maha Esa karena atas berkat dan anugerah-Nya sehingga penulis dapat menyelesaikan Tugas Akhir yang berjudul:
\begin{center}
    \makeatletter
    ``\noindent\MakeUppercase{\@judul}''
    \makeatother
\end{center}

\noindent \Ta~ini disusun guna memenuhi salah satu syarat kelulusan pada jenjang Sarjana di Departemen Matematika, Fakultas Sains dan Analitika Data, Institut Teknologi Sepuluh Nopember. \Ta~ini tentu tidak akan selesai tanpa dukungan, kerja sama, dan bantuan dari banyak pihak. Oleh karena itu, penulis ingin menyampaikan rasa terima kasih yang mendalam kepada:
%
\begin{enumerate}
    \item Kedua orang tua penulis yang senantiasa memberikan doa, dukungan, dan dorongan selama perkuliahan.
    \item Bapak Dr. Didik Khusnul Arif, S.Si., M.Si., selaku Kepala Departemen Matematika FSAD ITS atas dukungan sarana prasarana yang diberikan kepada penulis.
    \item Bapak Dr. Drs. Bandung Arry S., MI.Komp., selaku dosen pembimbing yang selalu memberikan motivasi, arahan, dukungan moral dan fisik selama penyusunan \ta.
    \item Bapak/Ibu dosen pengajar yang tidak dapat disebutkan satu per satu atas ilmu yang diberikan, serta seluruh staf, tenaga kependidikan, dan keluarga besar Departemen Matematika, Fakultas Sains dan Analitika Data, ITS, atas bantuan dan dukungan selama perkuliahan berlangsung.
    \item Teman-teman seperjuangan yang telah senantiasa menemani serta memberikan dukungan dan bantuan selama perkuliahan.
\end{enumerate}

Penulis menyadari bahwa \ta~masih jauh dari kata sempurna. Oleh karena itu, segala bentuk masukan maupun saran yang membangun dari para pembaca sangat diharapkan guna untuk pengembangan penulisan kedepannya. Akhir kata, sekiranya \ta~ dapat memberikan manfaat bagi semua pihak yang berkepentingan.

\makeatletter
\begin{flushright}
    Surabaya, \ \@tanggal\ \@bulan\ \@tahun

    \vspace{2cm}
    \@nama
\end{flushright}
\makeatother